% Libro 2 · Mundo 6 · Secuencias y patrones (app: pu-secuencias)
\mundo{Secuencias y patrones}{pu-secuencias}{¿Qué número sigue? ¿Qué letra falta? ¿Qué figura ocupa el lugar 100? Detrás
de cada secuencia hay una regla, y el trabajo es descubrirla con método: mirar diferencias, cocientes, posiciones y
ciclos, en ese orden.}{%
\item \textbf{6.1} Sucesiones numéricas
\item \textbf{6.2} Letras y figuras
\item \textbf{6.3} Patrones cíclicos}

\leccion{pu-se-1}{Sucesiones numéricas}{Juego del término que sigue.}

\begin{reto}
¿Qué número sigue en la sucesión $2,\ 5,\ 11,\ 23,\ 47,\ \ldots$?
\opciones{94}{95}{96}{71}
\end{reto}

\subsection{La idea: preguntarle a la sucesión cómo cambia}
El método es una lista de preguntas, en orden:
\begin{enumerate}
  \item ¿Las \clave{diferencias} son constantes? Es aritmética: $3, 7, 11, 15$ (suma 4).
  \item ¿Las diferencias forman su propia sucesión? Se calculan las \clave{segundas diferencias}.
  \item ¿Los \clave{cocientes} son constantes? Es geométrica: $3, 6, 12, 24$ (por 2).
  \item ¿Hay una mezcla: «por 2 y más 1», «por 3 y menos 2»?
  \item ¿Son dos sucesiones \clave{intercaladas} (los lugares pares siguen una regla y los impares otra)?
  \item ¿Son números conocidos: cuadrados, cubos, primos, Fibonacci (cada término es la suma de los dos anteriores)?
\end{enumerate}

\begin{center}
\begin{tikzpicture}[x=11mm,font=\small]
  \foreach \i/\v in {0/1,1/4,2/9,3/16,4/25}{\node at (\i,0) {$\v$};}
  \foreach \i/\v in {0/3,1/5,2/7,3/9}{\node[text=zolive,font=\scriptsize] at (\i+.5,-.55) {$+\v$};}
  \foreach \i in {0,1,2}{\node[text=zcoral,font=\scriptsize] at (\i+1,-1.05) {$+2$};}
  \node[right,font=\scriptsize,align=left] at (5,-.5) {segundas diferencias constantes:\\una sucesión cuadrática};
\end{tikzpicture}
\end{center}

\subsubsection{El término general}
En una aritmética, $a_n = a_1 + (n - 1)d$: el término 50 de $3, 7, 11, \ldots$ es $3 + 49 \cdot 4 = 199$. En una
geométrica, $a_n = a_1 r^{n-1}$.

\begin{regla}
\begin{itemize}
  \item Aritmética: $a_n = a_1 + (n - 1)d$; suma de $n$ términos: $\tfrac{n(a_1 + a_n)}{2}$.
  \item Geométrica: $a_n = a_1 r^{n - 1}$.
  \item Orden de búsqueda: diferencias, segundas diferencias, cocientes, mezclas, intercaladas, números conocidos.
\end{itemize}
\end{regla}

\subsection{Ejemplos resueltos}
\begin{ejemplo}[Segundas diferencias]
$3,\ 4,\ 7,\ 12,\ 19,\ \ldots$
\solucion Diferencias: 1, 3, 5, 7 (impares). Sigue $+9$: 28.
\end{ejemplo}

\begin{ejemplo}[Intercaladas]
$1,\ 20,\ 3,\ 17,\ 9,\ 14,\ 27,\ \ldots$
\solucion Lugares impares: 1, 3, 9, 27 ($\times 3$). Pares: 20, 17, 14 ($-3$). Sigue el lugar 8 (par): 11.
\end{ejemplo}

\begin{ejemplo}[Suma de Gauss]
¿Cuánto suman los números del 1 al 200?
\solucion $\tfrac{200 \cdot 201}{2} = 20\,100$. (Emparejando el primero con el último: 100 parejas que suman 201.)
\end{ejemplo}

\begin{trampa}
Quedarse con la primera regla que parece funcionar sin comprobarla con \emph{todos} los términos dados.
\end{trampa}

\taller
\piso{Entrena}{reglas}
\begin{ejercicios}[2]
\item $5, 9, 13, 17, \ldots$ \sol{21.}
\item $2, 6, 18, 54, \ldots$ \sol{162.}
\item $1, 1, 2, 3, 5, 8, \ldots$ \sol{13.}
\item $2, 3, 5, 9, 17, \ldots$ \sol{33 ($\times 2 - 1$).}
\item $1, 8, 27, 64, \ldots$ \sol{125.}
\item $100, 98, 94, 88, \ldots$ \sol{80 ($-2, -4, -6, -8$).}
\end{ejercicios}
\piso{Como en el examen}{opción múltiple}
\begin{ejercicios}
\item El término 40 de la sucesión $7, 12, 17, 22, \ldots$ es
\opciones{202}{207}{197}{200}
\sol{A. $7 + 39 \cdot 5$.}
\item ¿Qué número falta? $\;3,\ 7,\ 15,\ \square,\ 63$
\opciones{29}{31}{30}{33}
\sol{B. $\times 2 + 1$.}
\end{ejercicios}

\begin{retoresuelto}
Cada término es el doble del anterior más 1: $2 \to 5 \to 11 \to 23 \to 47 \to 95$. (También: las diferencias 3, 6, 12,
24 se duplican; sigue $+48$.) \textbf{Opción B.}
\end{retoresuelto}
\bitacora

\leccion{pu-se-2}{Letras y figuras}{Juego del alfabeto numerado y las figuras que giran.}

\begin{reto}
¿Qué letra sigue en la sucesión A, C, F, J, O, …? (Usa el alfabeto sin Ñ.)
\opciones{T}{U}{V}{S}
\end{reto}

\subsection{La idea: convertir letras en números}
Las letras se vuelven números con su posición en el alfabeto. El examen suele usar el alfabeto de 26 letras (sin Ñ, CH ni
LL); si usa otro, lo dice. Conviene tener la tabla a mano:

\begin{center}\scriptsize\setlength{\tabcolsep}{2.3pt}
\begin{tabular}{*{13}{c}}
A & B & C & D & E & F & G & H & I & J & K & L & M\\
1 & 2 & 3 & 4 & 5 & 6 & 7 & 8 & 9 & 10 & 11 & 12 & 13\\[3pt]
N & O & P & Q & R & S & T & U & V & W & X & Y & Z\\
14 & 15 & 16 & 17 & 18 & 19 & 20 & 21 & 22 & 23 & 24 & 25 & 26\\
\end{tabular}
\end{center}

\subsubsection{Figuras}
En las sucesiones de figuras, cada característica cambia por su cuenta: la \emph{posición} (gira $90^\circ$ cada vez),
el \emph{número} de elementos (aumenta de a uno), el \emph{color} (alterna), el \emph{tamaño}. Se estudia una
característica a la vez.

\begin{center}
\begin{tikzpicture}[x=12mm]
  \foreach \i/\a in {0/0,1/90,2/180,3/270}{%
    \draw[zink,line width=.6pt] (\i,0) circle (4mm);
    \fill[zolive] (\i,0) ++(\a:2.4mm) circle (1.2mm);}
  \node at (4,0) {?};
  \node[right,font=\scriptsize] at (4.4,0) {el punto gira $90^\circ$: vuelve a la derecha};
\end{tikzpicture}
\end{center}

\begin{regla}
\begin{itemize}
  \item Letras a números con su posición; luego, el método de las sucesiones numéricas.
  \item Si se pasa de la Z, se vuelve a empezar (después de Z viene A).
  \item En figuras, separa las características y estudia cada una.
\end{itemize}
\end{regla}

\subsection{Ejemplos resueltos}
\begin{ejemplo}[Pares de letras]
AZ, BY, CX, DW, …
\solucion La primera avanza (A, B, C, D, E) y la segunda retrocede (Z, Y, X, W, V): EV. Además, cada par suma 27.
\end{ejemplo}

\begin{ejemplo}[Saltos crecientes]
B, E, J, Q, …
\solucion 2, 5, 10, 17: diferencias 3, 5, 7; sigue $+9$: 26, Z.
\end{ejemplo}

\begin{trampa}
Usar un alfabeto con Ñ cuando el examen usa el de 26 letras (o al revés): todo se corre una posición después de la N.
\end{trampa}

\taller
\piso{Entrena}{letras}
\begin{ejercicios}[2]
\item C, F, I, L, … \sol{O.}
\item Z, X, U, Q, … \sol{L ($-2, -3, -4, -5$).}
\item A, B, D, H, … \sol{P (1, 2, 4, 8, 16).}
\item AB, DE, GH, … \sol{JK.}
\end{ejercicios}
\piso{Como en el examen}{opción múltiple}
\begin{ejercicios}
\item ¿Qué par sigue? $\;$ BD, FH, JL, …
\opciones{NP}{MO}{NO}{OQ}
\sol{A. Saltos de 2 dentro del par y de 4 entre pares.}
\end{ejercicios}

\begin{retoresuelto}
Posiciones: 1, 3, 6, 10, 15 (diferencias 2, 3, 4, 5). Sigue $15 + 6 = 21$: U. \textbf{Opción B.}
\end{retoresuelto}
\bitacora

\leccion{pu-se-3}{Patrones cíclicos}{Juego de la rueda que da vueltas.}

\begin{reto}
Un semáforo está en verde 40 segundos, en amarillo 5 y en rojo 35, y repite. Si se pone en verde a las 8:00:00, ¿qué
color tiene a las 8:20:00?
\opciones{verde}{amarillo}{rojo}{no se puede saber}
\end{reto}

\subsection{La idea: el residuo dice dónde estás}
Todo lo que se repite con un período fijo (días de la semana, colores, letras de una palabra, cifras de un decimal
periódico) se resuelve con el residuo: si el ciclo mide $k$, la posición $N$ es igual a la posición $r$, donde $r$ es el
residuo de $N \div k$.

\begin{center}
\begin{tikzpicture}[font=\scriptsize]
  \draw[zline,line width=4pt] (0,0) circle (1);
  \draw[zlime!80!black,line width=4pt] (90:1) arc[start angle=90,end angle=-90,radius=1];
  \draw[zsun,line width=4pt] (-90:1) arc[start angle=-90,end angle=-112.5,radius=1];
  \draw[zcoral,line width=4pt] (-112.5:1) arc[start angle=-112.5,end angle=-270,radius=1];
  \node at (0,0) {80 s};
  \node[right] at (1.3,.3) {verde: 0 a 40};
  \node[right] at (1.3,-.1) {amarillo: 40 a 45};
  \node[right] at (1.3,-.5) {rojo: 45 a 80};
\end{tikzpicture}
\end{center}

\begin{regla}
\begin{itemize}
  \item Mide el ciclo $k$. Calcula $r = N \bmod k$.
  \item Ojo con dónde empieza la cuenta: si el primer elemento es la posición 1, el residuo 0 corresponde al último del
  ciclo.
  \item Si hay dos ciclos a la vez, se repiten juntos cada MCM de sus longitudes.
\end{itemize}
\end{regla}

\subsection{Ejemplos resueltos}
\begin{ejemplo}[Palabra que se repite]
Se escribe ZUKAZUKAZUKA… ¿Cuántas A hay entre las primeras 2\,026 letras?
\solucion Bloques de 4 con una A cada uno: $2\,026 = 506 \cdot 4 + 2$. Hay 506 bloques completos (506 A) y dos letras
más (Z, U). Total 506.
\end{ejemplo}

\begin{ejemplo}[Dos ciclos]
Una luz roja parpadea cada 4 segundos y una azul cada 6. Si parpadean juntas en $t = 0$, ¿cuántas veces coinciden en
el primer minuto (incluido $t = 0$)?
\solucion Coinciden cada MCM$(4, 6) = 12$ s: en 0, 12, 24, 36, 48 y 60: 6 veces.
\end{ejemplo}

\begin{trampa}
Contar desde 0 cuando el problema cuenta desde 1 (o al revés). Haz la prueba con un caso pequeño antes de usar el
residuo.
\end{trampa}

\taller
\piso{Entrena}{ciclos}
\begin{ejercicios}
\item Si el 1 de marzo es lunes, ¿qué día es el 1 de abril? \sol{Marzo tiene 31 días: $31 \bmod 7 = 3$: jueves.}
\item ¿Qué figura ocupa el lugar 50 en ○ △ □ ○ △ □ …? \sol{$50 \bmod 3 = 2$: △.}
\item Un faro gira y alumbra la playa cada 12 s. ¿Cuántas veces la alumbra en una hora? \sol{$3\,600 / 12 = 300$.}
\end{ejercicios}
\piso{Como en el examen}{opción múltiple}
\begin{ejercicios}
\item En la secuencia 1, 2, 3, 4, 3, 2, 1, 2, 3, 4, 3, 2, 1, … ¿qué número ocupa el lugar 100?
\opciones{1}{2}{3}{4}
\sol{D. Ciclo 1, 2, 3, 4, 3, 2 de longitud 6. $100 = 16 \cdot 6 + 4$: el cuarto, 4.}
\end{ejercicios}

\begin{retoresuelto}
De las 8:00:00 a las 8:20:00 hay 1\,200 s. El ciclo mide $40 + 5 + 35 = 80$ s, y $1\,200 = 15 \cdot 80$: se completan
exactamente 15 ciclos y el semáforo acaba de ponerse en verde. \textbf{Opción A.}
\end{retoresuelto}
\bitacora

\evaluacion
\begin{ejercicios}
\item ¿Qué número sigue? $\;1,\ 2,\ 6,\ 24,\ 120,\ \ldots$
\opciones{240}{600}{720}{144}
\sol{C. Factoriales: $\times 2, \times 3, \ldots, \times 6$. (6.1)}
\item El término 25 de $4, 11, 18, 25, \ldots$ es
\opciones{172}{175}{168}{179}
\sol{A. $4 + 24 \cdot 7$. (6.1)}
\item ¿Qué letra sigue? $\;$ D, G, J, M, …
\opciones{O}{P}{Q}{N}
\sol{B. Saltos de 3. (6.2)}
\item ¿Qué par sigue? $\;$ AY, CW, EU, …
\opciones{GS}{FT}{GT}{HS}
\sol{A. Primera $+2$, segunda $-2$. (6.2)}
\item Si hoy es domingo, dentro de 75 días será
\opciones{jueves}{viernes}{miércoles}{sábado}
\sol{B. $75 \bmod 7 = 5$. (6.3)}
\item En ABCDABCDABCD…, la letra 2\,027 es
\opciones{A}{B}{C}{D}
\sol{C. $2\,027 \bmod 4 = 3$. (6.3)}
\end{ejercicios}
